Across the twenty matched fleet-size/radius comparisons, the 20 m/s policy has
the smaller mean time-averaged distance in all 20 cases. The median paired
difference $(d_{12}-d_{20})$ is +51.8 km. Uncorrected two-sided paired tests
indicate lower distance at 20 m/s in 19 cases and at 12 m/s in none of the
cases at $p<0.05$. Because these comparisons were not adjusted for
multiplicity, they provide exploratory rather than confirmatory evidence.

Averaged over the four fleet sizes, the differences (12 m/s minus 20 m/s) at
the five radii are: 100~km: $\Delta d=+10.1$~km, $\Delta L=-78.7$~km,
$\Delta \bar v=-1.74$~m/s; 200~km: $\Delta d=+19.3$~km,
$\Delta L=-95.5$~km, $\Delta \bar v=-2.07$~m/s; 400~km:
$\Delta d=+49.1$~km, $\Delta L=-147.4$~km, $\Delta \bar v=-3.02$~m/s;
600~km: $\Delta d=+73.5$~km, $\Delta L=-196.6$~km,
$\Delta \bar v=-3.79$~m/s; 800~km: $\Delta d=+84.2$~km,
$\Delta L=-233.3$~km, $\Delta \bar v=-4.18$~m/s. Here, positive distance and
negative path-length differences mean that the 12 m/s policy is farther from
the center and travels less; negative mean-speed differences reflect its lower
propulsion use.

The relative time-averaged distance penalty of the 12 m/s policy is 100~km:
+6.3\%; 200~km: +10.6\%; 400~km: +18.2\%; 600~km: +17.7\%; 800~km: +14.0\%.
The absolute difference therefore increases with deployment radius, although
not monotonically in relative terms. This pattern is consistent with the finite
time available to remove forecast-induced displacement. The shorter 12 m/s
paths reduce propulsion demand, but the reduced speed is associated with
weaker TC-center tracking in the tested episodes.
